% =============================================================================
% Project : THE_PHILOSOPHY_OF_VIABILITY
% Book    : Book I — Ontology of Viability
% Book ID : VS.B01
%
% File    : 4_ASSETS/tables/
%           VS.B01.P01.C02.S09.SS02.P004.TL01.analytical_forms_of_invariance.tex
%
% Asset   : TL01 — Analytical Forms of Invariance
% Location: Chapter 02 / Section 09 / Subsection 02 / Paragraph 004
%
% Purpose :
%   To distinguish the principal analytical forms of invariance
%   by identifying their respective objects of preservation
%   and limits of inference.
%
% =============================================================================

\begin{table}[htbp]
    \centering
    \small
    \renewcommand{\arraystretch}{1.35}
    \setlength{\tabcolsep}{5pt}

    \begin{tabular}{|p{3.0cm}|p{5.2cm}|p{5.2cm}|}
        \hline
        \textbf{Form} &
        \textbf{Determination Examined} &
        \textbf{Limit of Inference}
        \\
        \hline

        Determinational &
        A specified determination or feature. &
        Does not establish the preservation of the
        existent as a whole or its numerical identity.
        \\
        \hline

        Structural &
        An organizational pattern, arrangement,
        or structural constraint. &
        Does not establish the preservation of
        every component or numerical identity.
        \\
        \hline

        Relational &
        A specified relation between relata. &
        Does not establish that every relatum
        remains unchanged or numerically identical.
        \\
        \hline

        Compositional &
        A specified determination concerning
        the composition of an existent. &
        Does not establish the preservation of
        every component or the whole existent's identity.
        \\
        \hline

        Quantitative &
        A specified magnitude, number, ratio,
        or other quantitative determination. &
        Does not establish qualitative, structural,
        or ontological preservation beyond its criterion.
        \\
        \hline

        Qualitative &
        A specified qualitative determination. &
        Does not establish quantitative or
        numerical identity.
        \\
        \hline

        Configurational &
        A specified configuration or arrangement
        of determinations. &
        Does not establish the preservation of
        every determination or component.
        \\
        \hline

        Relational-Structural &
        An organized pattern of relations under
        a specified organizational criterion. &
        Does not establish that every individual
        relation or relatum remains unchanged.
        \\
        \hline
    \end{tabular}

    \caption{Analytical forms of invariance and the respective
    determinations they examine. The listed forms are not
    mutually exclusive and do not constitute an exhaustive
    ontological taxonomy.}
    \label{tab:analytical-forms-of-invariance}
\end{table}

% =============================================================================
% END OF TABLE TL01
% =============================================================================